\documentclass[11pt]{article}

% ===================== PACKAGES =====================
\usepackage[a4paper,margin=1in]{geometry}
\usepackage{amsmath,amssymb}
\usepackage{enumitem}
\usepackage{fancyhdr}
\usepackage{xcolor} 

% ===================== HEADER & FOOTER =====================
\pagestyle{fancy}
\fancyhf{}
\lhead{CSE 400: Fundamentals of Probability in Computing}
\rhead{Lecture-7 Scribe Submission}
\cfoot{\thepage}

% ===================== CUSTOM COMMANDS =====================
\newcommand{\solution}{
    \vspace{0.5cm}
    \noindent\textbf{\textcolor{blue}{Solution:}}
    \vspace{0.2cm}
    
}

% ===================== TITLE =====================
\title{
    \normalsize School of Engineering and Applied Science (SEAS), Ahmedabad University \\
    \vspace{0.2cm}
    \textbf{CSE 400: Fundamentals of Probability in Computing}\\
    \Large Lecture-7 Scribe Submission
}
\author{} 
\date{}

\begin{document}
\maketitle

\vspace{-2cm}
\begin{center}
    \begin{tabular}{ll}
        \textbf{Name:} & {Jashvi Shah \hspace{2.5in}} \\ [1.5ex]
        \textbf{Enrollment No:} & {AU2440026 \hspace{2.5in}} \\ [1.5ex]
        \textbf{Email:} & {jashvi.s@ahduni.edu.in \hspace{2.1in}} \\ [1.5ex]
        \textbf{Date of Submission:} & {Febrauary 7\textsuperscript{th}, 2026 (11:59 PM) \hspace{1.2in}}
    \end{tabular}
\end{center}

\hrule
\vspace{0.5cm}

\section*{CSE400 – Lecture 7 Scribe}


\textbf{Lecture:} Expectation, CDFs, PDFs and Problem Solving\\
\textbf{Instructor:} Dhaval Patel\\
\textbf{Date:} Jan 27, 2025

\bigskip

\section*{1. Lecture Structure}

According to the outline on page 2, the lecture covers:

\begin{itemize}
\item Cumulative Distribution Function (CDF)
\begin{itemize}
\item Definition
\item Properties
\item Example
\end{itemize}
\item Probability Density Function (PDF)
\begin{itemize}
\item Definition
\item Properties
\item Example
\end{itemize}
\item Expectation of Random Variables
\begin{itemize}
\item Definition
\item Expectation of function of RV
\item Linearity
\item Moments \& Central Moments
\item Variance
\item Skewness
\item Kurtosis
\end{itemize}
\item L7\_S2\_A\_Revised
\end{itemize}

\section*{2. Cumulative Distribution Function (CDF)}

\subsection*{2.1 Conceptual Illustration}

The diagram on pages 3–4 shows a cylindrical tank analogy:

\begin{itemize}
\item Height corresponds to value of random variable
\item Accumulated volume corresponds to cumulative probability
\item Maximum filled volume corresponds to probability 1
\end{itemize}

This visual introduces cumulative accumulation behavior.

\subsection*{2.2 Definition}

For random variable $X$:

\[
F_X(x)=Pr(X \le x)
\]

This function characterizes the probability behavior of the experiment.

\subsection*{2.3 Properties}

Bounds
\[
0 \le F_X(x) \le 1
\]

Limits
\[
F_X(-\infty)=0
\]
\[
F_X(\infty)=1
\]

Monotonicity
\[
x_1 < x_2 \Rightarrow F_X(x_1)\le F_X(x_2)
\]

Interval Probability
\[
Pr(x_1<X\le x_2)=F_X(x_2)-F_X(x_1)
\]

\subsection*{2.4 Example (Validity)}

Slides present candidate functions and verify whether they satisfy CDF properties by checking:

\begin{itemize}
\item Range constraint
\item Monotonic increase
\item Proper limits
\end{itemize}

Functions violating these conditions are rejected.

\subsection*{2.5 Example (Probability Computation)}

Given
\[
F_X(x)=(1-e^{-x})u(x)
\]

Probabilities computed using:

\[
Pr(X>a)=1-F_X(a)
\]

\[
Pr(a<X<b)=F_X(b)-F_X(a)
\]

\[
Pr(X>a \mid X<b)=\frac{F_X(b)-F_X(a)}{F_X(b)}
\]

\section*{3. Probability Density Function (PDF)}

\subsection*{3.1 Definition}

\[
f_X(x)=\lim_{\epsilon\to 0}\frac{Pr(x\le X<x+\epsilon)}{\epsilon}
\]

Using relation with CDF:

\[
f_X(x)=\frac{d}{dx}F_X(x)
\]

\subsection*{3.2 Fundamental Relationship}

\[
F_X(x)=\int_{-\infty}^{x} f_X(y)\,dy
\]

Thus:

PDF = derivative of CDF

CDF = integral of PDF

\subsection*{3.3 Properties}

Non-negativity
\[
f_X(x)\ge 0
\]

Normalization
\[
\int_{-\infty}^{\infty} f_X(x)\,dx = 1
\]

Interval Probability
\[
Pr(a<X\le b)=\int_a^b f_X(x)\,dx
\]

\subsection*{3.4 Examples}

Slides show checking candidate PDFs by verifying:

\begin{itemize}
\item Non-negative over domain
\item Integrates to 1
\end{itemize}

Conversion examples:

From CDF $\rightarrow$ PDF

Differentiate given CDF expression.

From PDF $\rightarrow$ CDF

Integrate given PDF expression.

\section*{4. Expectation of Random Variables}

\subsection*{4.1 Definition}

Discrete case:

\[
E[X]=\sum x_k\, p_X(x_k)
\]

Continuous case:

\[
E[X]=\int x f_X(x)\,dx
\]

L7\_S2\_A\_Revised

\subsection*{4.2 Expectation of Function of RV}

\[
E[g(X)]=\sum g(x_k)p_X(x_k)
\]

or

\[
\int g(x)f_X(x)\,dx
\]

\subsection*{4.3 Linearity of Expectation}

\[
E[aX+b]=aE[X]+b
\]

\section*{5. Moments and Central Moments}

\subsection*{5.1 n-th Moment}

\[
E[X^n]
\]

\subsection*{5.2 Central Moment}

\[
E[(X-\mu)^n]
\]

with

\[
\mu=E[X]
\]

\subsection*{5.3 Variance}

\[
\sigma^2 = E[X^2]-\mu^2
\]

\subsection*{5.4 Higher-Order Shape Measures}

Slides identify:

\begin{itemize}
\item Skewness
\item Kurtosis
\end{itemize}

as derived from higher central moments.

\section*{6. Logical Flow Summary}

Lecture progression:

\begin{itemize}
\item Introduce cumulative probability representation
\item Define and analyze CDF
\item Derive PDF and relationships
\item Apply through examples
\item Introduce expectation and linearity
\item Extend to moments and distribution shape
\end{itemize}

\end{document}
